\section{From the Optimizer History to Its Stationary Limit}
\label{app:transfer}

\paragraph{Keeping the optimizer state bounded.}
A stationary, or invariant, law is a probability distribution unchanged
by one optimizer step. The native cone is the set of moving-average
states satisfying the inequality below. For any scalar history, weighted
Cauchy--Schwarz gives
$m^2\leq C_*^2v$. Updating one sample preserves this relation: if
$m^2\leq C_*^2v$, weighted Cauchy--Schwarz gives
\[
 [\alpha m+(1-\alpha)g]^2
 \leq\left(\frac{\alpha^2C_*^2}{\beta}
           +\frac{(1-\alpha)^2}{1-\beta}\right)
       [\beta v+(1-\beta)g^2]
 =C_*^2[\beta v+(1-\beta)g^2].
\]
The normalized coordinate action has magnitude at most $C_*$, including
its continuous zero value at $m=v=0$. Hence, for $\eta\lambda\leq1$,
\[
 |u_{i,n+1}|\leq(1-\eta\lambda)|u_{i,n}|+\eta C_*.
\]
The box $|u_i|\leq M=C_*/\lambda$ is invariant. Every stationary cone law
is supported there: if $0<\eta\lambda<1$, for $r>0$ its tail above $M+r$ is bounded by its tail
above $M+r/(1-\eta\lambda)$, and iteration sends the latter threshold to
infinity. At $\eta\lambda=1$ confinement holds after one step.
The same shrinking-bound argument keeps the stationary moving averages
in a compact set because gradients are bounded on this box. The update
is continuous on a compact set that it preserves. It therefore has the
Feller property: the one-step expectation of a continuous function
is continuous in the current state.
Averaging successive laws gives an invariant probability. All bounds
hold uniformly over every such probability.

\paragraph{How fast the update forgets old samples.}
Write $r=\|x\|_d$ and let $a$ represent the numerator functional in
the weighted square-sum space of histories, so $\|a\|_d=C_*$. For $r>0$, with $z=x/r$,
\[
 DT(x)=\frac{[I-\tfrac r{r+\epsilon}(z\otimes z)]a}{r+\epsilon}.
\]
It extends continuously by $a/\epsilon$ at zero. Differentiating away
from zero gives $\|D^2T\|\leq6C_*/\epsilon^2$: the terms with
$(a\otimes z+z\otimes a)$ and $\langle a,z\rangle I$ have total
norm at most $3C_*/(r+\epsilon)^2$, and the remaining rank-one term
has norm at most $C_*(3r+\epsilon)/(r+\epsilon)^3$.
Integrating along a segment, with continuity at zero if the segment
crosses it, proves
$\|DT(x)-DT(y)\|\leq6C_*\|x-y\|_d/\epsilon^2$.
In particular $T$ is $C^1$ with bounded Lipschitz derivative.

For each coordinate, a frozen gradient history is affine in $u$ with bounded
parameter derivative. Two bounded mark histories sharing their last $k$
future samples differ in their weighted histories and derivatives by at
most $C\rho^k$. The update and its first parameter derivative therefore forget old
samples at a geometric rate. Let $\mathsf P$ be the history
shift with a fresh iid mark. Coupling to an independent stationary old
history proves uniform convergence of
\[
 H_u=\sum_{k\geq0}\mathsf P^k(f_u-F(u)),\qquad K_u=\mathsf P H_u,
\]
including its first parameter derivative on an enlarged box. Thus $H_u,K_u$
are bounded and uniformly Lipschitz in $u$. They give a Poisson corrector:
a bounded adjustment that separates predictable history effects from
fresh mean-zero noise. Only the update and its derivative need to forget
old samples; no total-variation mixing of the full infinite history is
required.

\paragraph{Using samples in the right order.}
In a stationary chain extended to past and future times, the contribution
of initial moving-average values vanishes.
The native increment bound gives $\|u_{n-j}-u_n\|\leq C\eta j$.
Since $\sum_jj^2d_j<\infty$, the Hilbert quotient Lipschitz bound gives
\[
 U_{n+1}^{\rm actual}=f_{u_n}(Z_{n+1})+O(\eta)
\]
uniformly. Here $Z_n$ is the complete past iid mark history, with the
current parameter measurable before the fresh mark. Define
\[
 \Delta M_{n+1}=H_{u_n}(Z_{n+1})-K_{u_n}(Z_n).
\]
It is bounded and has conditional mean zero given all past information,
including the stationary initial coordinate. Such increments are called
martingale differences. The Poisson identity is
\[
 f_{u_n}(Z_{n+1})-F(u_n)
 =\Delta M_{n+1}+K_{u_n}(Z_n)-K_{u_n}(Z_{n+1}).
\]
Consequently the negative correction $X_n=u_n-\eta K_{u_n}(Z_n)$ obeys
\begin{equation}\label{eq:corrected-core}
 X_{n+1}=X_n-\eta G(X_n)-\eta\Delta M_{n+1}+R_{n+1},
 \qquad \|R_{n+1}\|\leq C\eta^2.
\end{equation}
The $O(\eta^2)$ bound uses the parameter Lipschitz bounds for $K$ and $G$.
No restoring property of $G$ is needed for this representation. Its
conditional bracket $q(u,Z)$ is the conditional covariance matrix of the
next martingale difference. It is bounded and uniformly Lipschitz in $u$,
because $H_u$ is bounded and Lipschitz.

\begin{lemma}[Regular stationary transfer with a global Lyapunov metric]
\label{lem:regular-transfer}
A Lyapunov metric is a weighted squared distance that decreases under
the averaged dynamics. Suppose $G(0)=0$ and on an enlarged native box there
are a positive-definite
matrix $P_0$ and $c_0>0$ such that
\[
 u^\top P_0G(u)\geq c_0u^\top P_0u.
\]
Let $B=DG(0)$, and let $Q_0$ be the origin Green--Kubo matrix: the sum
of update covariances over all time lags at zero parameters. Every
invariant native-cone selection then satisfies
\[
 u/\sqrt\eta\Rightarrow N(0,S),\qquad BS+SB^\top=Q_0,
\]
with every positive absolute scaled moment converging. The corresponding
stationary path limit on clock $1/\eta$ is the stationary Ornstein--Uhlenbeck
(OU) diffusion: a linear diffusion with a Gaussian stationary law.
\end{lemma}
\begin{proof}
Use the norm $\|x\|_{P_0}^2=x^\top P_0x$. For each integer $k\geq1$,
stationary expansion of $\|X\|_{P_0}^{2k}$ in
\eqref{eq:corrected-core} gives a leading drift at most
$-2kc_0\eta\E\|X\|_{P_0}^{2k}$. The first martingale contribution
vanishes. Since corrected jumps are $O(\eta)$, the second-order remainder
is bounded by
\[
 C_k\eta^2\E\|X\|_{P_0}^{2k-2}+C_k\eta^{2k}.
\]
The remainder $R$ contributes at most
$C_k\eta^2\E\|X\|_{P_0}^{2k-1}$, bounded by the preceding even moment
on the fixed compact enlarged box. Thus
\[
 \E\|X\|_{P_0}^{2k}\leq C_k\eta
    (\E\|X\|_{P_0}^{2k-2}+\eta^{2k-2}).
\]
Induction and $\|u-X\|\leq C\eta$ prove
$\E\|u\|^{2k}=O(\eta^k)$ uniformly over invariant selections, before
any weak limit.

For the process $X_{\lfloor t/\eta\rfloor}/\sqrt\eta$, the bracket
parameter replacement costs, on every fixed time interval, at most
$C\eta\sum_{n\leq T/\eta}\E\|u_n\|=O(\sqrt\eta)$.
The remaining bracket at zero is a bounded function of the stationary iid
history. Its time average converges to its mean $Q_0$. The same convergence
holds uniformly over partial time sums by separating a fixed initial
segment. This uses only the history distribution; the initial parameter
may still be correlated with old samples.

The scaled martingale jumps are $O(\sqrt\eta)$, so no large jump survives
in the limit (the Lindeberg condition).
The scaled accumulated $R$ is $O(\sqrt\eta)$, as are corrector endpoints.
The moment bounds keep initial distributions from escaping to infinity.
Local Lipschitz drift, Doob's martingale inequality and Gronwall's inequality
then keep paths in a bounded region with high probability on every fixed
interval.
Differentiability at zero identifies the drift $-By$. Keeping the initial coordinate among the information used to define the
martingale shows that every subsequential stationary process limit solves
\[
 dY=-BYdt+Q_0^{1/2}dW.
\]
Differentiating the Lyapunov inequality at zero gives
$\operatorname{Sym}(P_0B)\succeq c_0P_0$, hence every eigenvalue of $B$
has positive real part. The OU diffusion has a unique invariant Gaussian
law, whose covariance is the unique solution of the displayed Lyapunov
equation. All subsequences therefore have that limit. The bounds on higher even moments give uniform integrability, allowing
expectations to pass to the limit and proving every claimed absolute
moment convergence.

Finally, its martingale origin variance equals the Green--Kubo matrix:
the bounded Poisson decomposition telescopes in each partial action sum;
dividing its covariance by the length and using absolute covariance
summability identifies the limit with $Q_0$. This also applies to every
scalar projection and hence to the matrix identity.
\qed\end{proof}
